% GENERATED FILE. Edit canonical lessons, then run npm run generate:books.
% canonical-source-hash: fnv1a64:51deb3813b72e55f
% canonical-lessons: PA-C183-habitual-return, PA-C183-karda, PA-C183-karde, PA-C183-utthda, PA-C183-dekhda, PA-C183-saunda, PA-C183-routine-review, PA-C183-payoff

\chapter{One Routine, One Familiar Engine}
\label{ch:pa-a2-habitual-routine}
\hlchaptermodality{183}

\begin{chapteropening}
\textbf{By the end of this chapter:} \emph{I can recognise and use five familiar verbs in exact present-habitual routine forms.}

Retrieve five exact present-habitual routine forms through listening, speaking, reading, and bounded copying after a slow review-led ramp.
\end{chapteropening}

\section[(habitual review)]{Return to one familiar engine}

\label{lesson:PA-C183-habitual-return}



\begin{quote}
Say the four familiar infinitives: \emph{karnaa}, \emph{dekhnaa}, \emph{utthnaa}, \emph{saunaa}.
\end{quote}

\begin{grammarlens}[title={Grammar Lens: stem plus a familiar ending}]
You already used \emph{-daa} for one masculine subject and \emph{-de} for a plural or respectful subject. Today we only bring that old engine back. No written target form appears yet.
\end{grammarlens}

\subsection*{Your turn}
Say these aloud:

\begin{itemize}
\raggedright
  \item remove the infinitive ending from \emph{karnaa}
  \item the Punjabi-shaped and Hindi-shaped twins for \emph{to see}
  \item \emph{to get up}, then \emph{to sleep}
\end{itemize}

\begin{tcolorbox}[breakable,colback=teal!4,colframe=teal!35!black,title={Before you move on}]
Which ending marks one masculine habitual form? (\emph{-daa}.) Which one can mark plural or respectful agreement? (\emph{-de}.)

This is retrieval of grammar and verbs already taught much earlier in the track.
\end{tcolorbox}
\section[\texorpdfstring{\pa{ਕਰਦਾ} (kardaa)}{ਕਰਦਾ (kardaa)}]{\pa{ਕਰਦਾ} — one man does or works}

\label{lesson:PA-C183-karda}



\begin{quote}
\emph{Hear:} \emph{kardaa} Hear the familiar stem \emph{kar-} and the familiar masculine ending \emph{-daa}.
\end{quote}

\begin{grammarlens}[title={Grammar Lens: the first exact form}]
\begin{quote}
\textbf{\pa{ਕਰਦਾ}} — \emph{kardaa} — does; works, for one masculine subject
\end{quote}

In a complete sentence you already know: \textbf{\pa{ਮੈਂ} \pa{ਕੰਮ} \pa{ਕਰਦਾ} \pa{ਹਾਂ}} — a man says, “I work.”
\end{grammarlens}

\subsection*{Your turn}
\begin{itemize}
\raggedright
  \item \emph{Hear:} \emph{kardaa}; YOU SAY: \textbf{\pa{ਕਰਦਾ}}
  \item \emph{Read it:} \textbf{\pa{ਕਰਦਾ}}
  \item \emph{Write it:} copy \textbf{\pa{ਕਰਦਾ}} once while saying \emph{kar-daa}
\end{itemize}

\begin{tcolorbox}[breakable,colback=teal!4,colframe=teal!35!black,title={Before you move on}]
Read the exact form for “does” with one masculine subject. (\textbf{\pa{ਕਰਦਾ}}.)

Source: the present-habitual engine and \emph{kamm karnaa} taught earlier in this track.
\end{tcolorbox}
\section[\texorpdfstring{\pa{ਕਰਦੇ} (karde)}{ਕਰਦੇ (karde)}]{\pa{ਕਰਦੇ} — more than one, or respectful agreement}

\label{lesson:PA-C183-karde}



\begin{quote}
Read yesterday’s one-man form: \textbf{\pa{ਕਰਦਾ}}.
\end{quote}

\begin{grammarlens}[title={Grammar Lens: change only the ending}]
\begin{quote}
\textbf{\pa{ਕਰਦੇ}} — \emph{karde} — do; work, with plural or respectful agreement
\end{quote}

The stem stays \emph{kar-}. The visible ending changes from \textbf{\pa{ਦਾ}} to \textbf{\pa{ਦੇ}}. In a complete sentence: \textbf{\pa{ਉਹ} \pa{ਕੰਮ} \pa{ਕਰਦੇ} \pa{ਹਨ}} — “They work,” or respectful “He works.”
\end{grammarlens}

\subsection*{Your turn}
\begin{itemize}
\raggedright
  \item \emph{Hear:} \emph{kardaa, karde}; YOU SAY: \textbf{\pa{ਕਰਦਾ}, \pa{ਕਰਦੇ}}
  \item \emph{Read it:} \textbf{\pa{ਕਰਦੇ}}
  \item \emph{Write it:} copy \textbf{\pa{ਕਰਦੇ}} once, changing only the ending
\end{itemize}

\begin{tcolorbox}[breakable,colback=teal!4,colframe=teal!35!black,title={Before you move on}]
Read the pair for one masculine subject, then plural or respectful agreement. (\textbf{\pa{ਕਰਦਾ}, \pa{ਕਰਦੇ}}.)

Source: the present-habitual agreement contrast taught earlier in this track.
\end{tcolorbox}
\section[\texorpdfstring{\pa{ਉੱਠਦਾ} (utthdaa)}{ਉੱਠਦਾ (utthdaa)}]{\pa{ਉੱਠਦਾ} — he gets up}

\label{lesson:PA-C183-utthda}



\begin{quote}
Recall \emph{utthnaa}, “to get up,” then \textbf{\pa{ਕਰਦਾ}}.
\end{quote}

\begin{grammarlens}[title={Grammar Lens: move the same ending}]
\begin{quote}
\textbf{\pa{ਉੱਠਦਾ}} — \emph{utthdaa} — gets up, for one masculine subject
\end{quote}

The infinitive was learned long ago. Keep its stem and use the same ending you saw in \textbf{\pa{ਕਰਦਾ}}. A tiny complete sentence is \textbf{\pa{ਉਹ} \pa{ਉੱਠਦਾ} \pa{ਹੈ}} — “He gets up.”
\end{grammarlens}

\subsection*{Your turn}
\begin{itemize}
\raggedright
  \item \emph{Hear:} \emph{utthdaa}; YOU SAY: \textbf{\pa{ਉੱਠਦਾ}}
  \item \emph{Read it:} \textbf{\pa{ਉੱਠਦਾ}}
  \item \emph{Write it:} copy \textbf{\pa{ਉੱਠਦਾ}} once, then point to the familiar \textbf{\pa{ਦਾ}} ending
\end{itemize}

\begin{tcolorbox}[breakable,colback=teal!4,colframe=teal!35!black,title={Before you move on}]
Say and read “he gets up,” then recall the two \emph{kar-} forms. (\textbf{\pa{ਉੱਠਦਾ}; \pa{ਕਰਦਾ}; \pa{ਕਰਦੇ}}.)

Source: the earlier \emph{utthnaa} lesson and the track's existing present-habitual engine.
\end{tcolorbox}
\section[\texorpdfstring{\pa{ਦੇਖਦਾ} (dekhdaa)}{ਦੇਖਦਾ (dekhdaa)}]{\pa{ਦੇਖਦਾ} — he sees or watches}

\label{lesson:PA-C183-dekhda}



\begin{quote}
Recall the twin infinitives \emph{vekhnaa} and \emph{dekhnaa}, then read \textbf{\pa{ਉੱਠਦਾ}}.
\end{quote}

\begin{grammarlens}[title={Grammar Lens: use the dekh- twin}]
\begin{quote}
\textbf{\pa{ਦੇਖਦਾ}} — \emph{dekhdaa} — sees or watches, for one masculine subject
\end{quote}

An earlier lesson taught that \emph{vekhnaa} and \emph{dekhnaa} are living twins. This mock-facing form uses the \emph{dekh-} twin: \textbf{\pa{ਉਹ} \pa{ਦੇਖਦਾ} \pa{ਹੈ}} — “He sees” or “He watches.”
\end{grammarlens}

\subsection*{Your turn}
\begin{itemize}
\raggedright
  \item \emph{Hear:} \emph{dekhdaa}; YOU SAY: \textbf{\pa{ਦੇਖਦਾ}}
  \item \emph{Read it:} \textbf{\pa{ਦੇਖਦਾ}}
  \item \emph{Write it:} copy \textbf{\pa{ਦੇਖਦਾ}} once while saying \emph{dekh-daa}
\end{itemize}

\begin{tcolorbox}[breakable,colback=teal!4,colframe=teal!35!black,title={Before you move on}]
Read “he sees,” then “he gets up.” (\textbf{\pa{ਦੇਖਦਾ}; \pa{ਉੱਠਦਾ}}.)

Source: the \emph{vekhnaa/dekhnaa} comparison taught earlier in this track.
\end{tcolorbox}
\section[\texorpdfstring{\pa{ਸੌਂਦਾ} (saundaa)}{ਸੌਂਦਾ (saundaa)}]{\pa{ਸੌਂਦਾ} — he sleeps}

\label{lesson:PA-C183-saunda}



\begin{quote}
Recall \emph{saunaa}, “to sleep,” then read \textbf{\pa{ਦੇਖਦਾ}}.
\end{quote}

\begin{grammarlens}[title={Grammar Lens: one small sound change}]
\begin{quote}
\textbf{\pa{ਸੌਂਦਾ}} — \emph{saundaa} — sleeps, for one masculine subject
\end{quote}

Hear the nasal in the middle and keep the familiar ending. In a complete sentence: \textbf{\pa{ਉਹ} \pa{ਸੌਂਦਾ} \pa{ਹੈ}} — “He sleeps.”
\end{grammarlens}

\subsection*{Your turn}
\begin{itemize}
\raggedright
  \item \emph{Hear:} \emph{saundaa}; YOU SAY: \textbf{\pa{ਸੌਂਦਾ}}
  \item \emph{Read it:} \textbf{\pa{ਸੌਂਦਾ}}
  \item \emph{Write it:} copy \textbf{\pa{ਸੌਂਦਾ}} once; check the nasal mark before \textbf{\pa{ਦਾ}}
\end{itemize}

\begin{tcolorbox}[breakable,colback=teal!4,colframe=teal!35!black,title={Before you move on}]
Read the three routine forms: \textbf{\pa{ਉੱਠਦਾ}, \pa{ਦੇਖਦਾ}, \pa{ਸੌਂਦਾ}}.

Source: the earlier \emph{saunaa} lesson and the track's existing present-habitual engine.
\end{tcolorbox}
\section[\texorpdfstring{\pa{ਕਰਦਾ}, \pa{ਕਰਦੇ}, \pa{ਉੱਠਦਾ} … (kardaa, karde …)}{ਕਰਦਾ, ਕਰਦੇ, ਉੱਠਦਾ … (kardaa, karde …)}]{Five forms, one engine}

\label{lesson:PA-C183-routine-review}



\begin{quote}
Read the agreement pair: \textbf{\pa{ਕਰਦਾ}, \pa{ਕਰਦੇ}}.
\end{quote}

\subsection*{Your turn: a short routine}
Read slowly, with the English meanings beside the line:

\begin{itemize}
\raggedright
  \item \textbf{\pa{ਉਹ} \pa{ਉੱਠਦਾ} \pa{ਹੈ}} — he gets up
  \item \textbf{\pa{ਉਹ} \pa{ਦੇਖਦਾ} \pa{ਹੈ}} — he watches
  \item \textbf{\pa{ਉਹ} \pa{ਸੌਂਦਾ} \pa{ਹੈ}} — he sleeps
  \item \textbf{\pa{ਮੈਂ} \pa{ਕੰਮ} \pa{ਕਰਦਾ} \pa{ਹਾਂ}} — I work, said by a man
  \item \textbf{\pa{ਉਹ} \pa{ਕੰਮ} \pa{ਕਰਦੇ} \pa{ਹਨ}} — they work, or respectful he works
\end{itemize}

\emph{Hear:} the five forms in mixed order; YOU SAY: each matching meaning

\begin{tcolorbox}[breakable,colback=teal!4,colframe=teal!35!black,title={Before you move on}]
\emph{Write it:} cover the model and write \textbf{\pa{ਕਰਦਾ}, \pa{ਕਰਦੇ}, \pa{ਉੱਠਦਾ}, \pa{ਦੇਖਦਾ}, \pa{ਸੌਂਦਾ}} once

Check the two agreement endings first; then check the three routine stems.

This review adds no new form and makes no mock-readiness claim.
\end{tcolorbox}
\section[\texorpdfstring{\pa{ਕਰਦਾ}, \pa{ਕਰਦੇ}, \pa{ਉੱਠਦਾ} … (kardaa, karde …)}{ਕਰਦਾ, ਕਰਦੇ, ਉੱਠਦਾ … (kardaa, karde …)}]{A gentle five-form payoff}

\label{lesson:PA-C183-payoff}



\begin{quote}
Without looking, say the \emph{kar-} pair for one masculine subject and for plural or respectful agreement.
\end{quote}

\subsection*{Your turn: four skills, bounded to five forms}
\begin{itemize}
\raggedright
  \item \emph{Hear:} \emph{utthdaa, dekhdaa, saundaa}; YOU SAY: \textbf{\pa{ਉੱਠਦਾ}, \pa{ਦੇਖਦਾ}, \pa{ਸੌਂਦਾ}}
  \item \emph{Read it:} \textbf{\pa{ਮੈਂ} \pa{ਕੰਮ} \pa{ਕਰਦਾ} \pa{ਹਾਂ।} \pa{ਉਹ} \pa{ਕੰਮ} \pa{ਕਰਦੇ} \pa{ਹਨ।}}
  \item \emph{Write it:} cover the model and write \textbf{\pa{ਕਰਦਾ}, \pa{ਕਰਦੇ}, \pa{ਉੱਠਦਾ}, \pa{ਦੇਖਦਾ}, \pa{ਸੌਂਦਾ}}
\end{itemize}

Now point to the ending that distinguishes \textbf{\pa{ਕਰਦਾ}} from \textbf{\pa{ਕਰਦੇ}}. Then say what stayed constant.

\begin{tcolorbox}[breakable,colback=teal!4,colframe=teal!35!black,title={Before you move on}]
Retrieve all five forms once more in any order.

This is the sixth prerequisite tranche exposed by Punjabi A2 mock drafting. It does not unblock Punjabi A2 mock 1 by itself. It is not an A2 mock, not A2 readiness or attainment, and not external or human validation.

Source: infinitives and the present-habitual agreement engine already taught earlier in the track.
\end{tcolorbox}
